\chapter{Taylor's Theorem and Approximation}

There is a quiet miracle hiding inside calculus, and this chapter is about finally
seeing it clearly. The miracle is this: if you know everything about a function at a
\emph{single point} --- its value, its slope, its curvature, the rate at which its
curvature changes, and so on through all the derivatives --- then in a great many
cases you know the function \emph{everywhere}. One point, infinitely interrogated,
reconstructs the whole curve. That is an astonishing claim, and like all astonishing
claims it deserves both a picture that makes it feel inevitable and a proof that
makes it honest.

We have been circling this idea for a while. The linear approximation
$f(x) \approx f(a) + f'(a)(x-a)$ used the value and the slope at $a$ to build the
best straight line through the point. It was good for $x$ near $a$ and got steadily
worse as you walked away. The natural question --- the one a curious person cannot
help asking --- is: \emph{why stop at a line?} If one derivative buys a line, what
does a second derivative buy? A parabola, surely, one that not only passes through
the point with the right slope but also bends the right way. And a third derivative
buys a cubic, and so on. Each new derivative should let the approximation hug the
curve a little longer before peeling away.

That is the Taylor polynomial. The dream that the polynomials might close in on the
function \emph{forever}, leaving no gap at all, is the Taylor \emph{series}. And the
precise accounting of exactly how big the gap is at any finite stage --- the thing
that turns a hopeful picture into a usable, error-bounded tool --- is \emph{Taylor's
theorem with remainder}. By the end of this chapter you will know not only how to
write these approximations down, but why they work, how far you can trust them, and
why the dream of perfect convergence sometimes comes true and sometimes, shockingly,
does not.

\section{From a line to a polynomial: matching derivatives}

Start with the linear approximation and stare at what made it good. Near a point $a$,
the line
\[
  L(x) = f(a) + f'(a)(x-a)
\]
agrees with $f$ in two respects: $L(a) = f(a)$ (same height) and $L'(a) = f'(a)$
(same slope). Those two matched conditions are the whole reason the line clings to
the curve near $a$. It is the unique line that gets both right.

\begin{intuitionbox}
Approximation by matching is a recurring trick in all of mathematics: to mimic a
complicated object near a point, build a simple object that agrees with it in as many
\emph{derivatives} as you can afford. A line matches two pieces of data (value and
slope). A parabola can match three (value, slope, and curvature). A degree-$n$
polynomial has $n+1$ free coefficients, so it can match $n+1$ pieces of data: the
value and the first $n$ derivatives. The more derivatives you pin down at $a$, the
longer the impostor passes for the real thing as you move away from $a$.
\end{intuitionbox}

So let us be greedy. We want a polynomial $P_n$ of degree $n$ such that
\[
  P_n(a) = f(a), \quad P_n'(a) = f'(a), \quad P_n''(a) = f''(a), \quad \dots, \quad
  P_n^{(n)}(a) = f^{(n)}(a).
\]
Match the function and its first $n$ derivatives at the single point $a$. The
question is what coefficients make this happen.

Here is the clean way to find them. Write the polynomial not in powers of $x$ but in
powers of $(x-a)$, because we only care about behavior near $a$:
\[
  P_n(x) = c_0 + c_1(x-a) + c_2(x-a)^2 + c_3(x-a)^3 + \cdots + c_n(x-a)^n.
\]
Now extract the coefficients one at a time by differentiating and then plugging in
$x=a$, which kills every term that still has a factor of $(x-a)$ in it. Watch the
pattern.

Plug in $x=a$ directly: every term but the first dies, leaving $P_n(a) = c_0$. So
$c_0 = f(a)$.

Differentiate once: $P_n'(x) = c_1 + 2c_2(x-a) + 3c_3(x-a)^2 + \cdots$. Plug in
$x=a$: everything but the first term dies, $P_n'(a) = c_1$. So $c_1 = f'(a)$.

Differentiate again: $P_n''(x) = 2c_2 + 3\cdot 2\, c_3 (x-a) + \cdots$. Plug in
$x=a$: $P_n''(a) = 2c_2$. So $c_2 = f''(a)/2$.

Once more: $P_n'''(x) = 3\cdot 2\, c_3 + 4\cdot 3\cdot 2\, c_4 (x-a) + \cdots$, so
$P_n'''(a) = 6c_3$, giving $c_3 = f'''(a)/6$.

The numbers $1, 2, 6, 24, \dots$ marching down the denominators are the factorials,
and they are not a coincidence: differentiating $(x-a)^k$ exactly $k$ times brings
down $k \cdot (k-1) \cdots 2 \cdot 1 = k!$ and turns the term into a constant, while
all lower terms have already vanished and all higher terms still carry a factor of
$(x-a)$. So in general
\[
  P_n^{(k)}(a) = k!\, c_k \quad\Longrightarrow\quad c_k = \frac{f^{(k)}(a)}{k!}.
\]

\begin{focusbox}
\textbf{The Taylor polynomial.} The unique degree-$n$ polynomial matching $f$ and its
first $n$ derivatives at $a$ is
\[
  P_n(x) = \sum_{k=0}^{n} \frac{f^{(k)}(a)}{k!}\,(x-a)^k
  = f(a) + f'(a)(x-a) + \frac{f''(a)}{2!}(x-a)^2 + \cdots + \frac{f^{(n)}(a)}{n!}(x-a)^n.
\]
The factorials are exactly what is needed to undo the factorials that
differentiation $k$ times produces. When $a=0$ this is sometimes called the
\emph{Maclaurin} polynomial, but it is the same idea centered at the origin.
\end{focusbox}

\begin{whybox}
Why factorials, really? Differentiation and the monomial $(x-a)^k$ are made for each
other. Each derivative peels one factor off the front: $(x-a)^k \to k(x-a)^{k-1} \to
k(k-1)(x-a)^{k-2} \to \cdots$. After exactly $k$ steps the variable is gone and a
bare $k!$ stands where the term used to be. Dividing the coefficient by $k!$ in
advance is the antidote: it guarantees that when you differentiate $k$ times and set
$x=a$, the $k$-th term contributes precisely $f^{(k)}(a)$ and nothing else
contributes at all. The factorial is not decoration; it is the gear ratio that makes
the $k$-th coefficient control the $k$-th derivative cleanly and independently of all
the others.
\end{whybox}

\subsection{A first picture: parabola beats line}

Take $f(x) = \cos x$ at $a = 0$. We know $f(0)=1$, $f'(0) = -\sin 0 = 0$,
$f''(0) = -\cos 0 = -1$. So
\[
  P_1(x) = 1, \qquad P_2(x) = 1 - \tfrac{1}{2}x^2.
\]
The linear approximation of cosine at the top of its hump is just the flat line
$y=1$ --- correct in value and slope but blind to the fact that cosine curves
\emph{downward}. The quadratic $1 - \tfrac{1}{2}x^2$ knows about that downward bend,
and you can see it follow the cosine far better.

\begin{figure}[H]
\centering
\begin{tikzpicture}
\begin{axis}[
  width=12cm, height=7cm,
  axis lines=middle,
  xlabel={$x$}, ylabel={$y$},
  xmin=-3.2, xmax=3.2, ymin=-1.6, ymax=1.5,
  samples=160, domain=-3.1:3.1,
  legend pos=south east, legend cell align=left,
  clip=true,
]
\addplot[curveblue, very thick] {cos(deg(x))};
\addlegendentry{$\cos x$}
\addplot[gray, thick, dashed, domain=-3.1:3.1] {1};
\addlegendentry{$P_1 = 1$}
\addplot[tangentred, thick] {1 - 0.5*x^2};
\addlegendentry{$P_2 = 1 - x^2/2$}
\addplot[only marks, mark=*, mark size=1.6pt, black] coordinates {(0,1)};
\end{axis}
\end{tikzpicture}
\caption{At $a=0$ the linear approximation to $\cos x$ is the flat line $y=1$. The
quadratic $1-\tfrac12 x^2$ captures the downward curvature and hugs the cosine over
a much wider stretch.}
\label{ch31:fig-cos-parabola}
\end{figure}

Each extra term tightens the grip. The degree-$4$ Taylor polynomial of cosine,
$1 - \tfrac12 x^2 + \tfrac{1}{24}x^4$, tracks the true cosine accurately well past
$x=2$. We will quantify ``accurately'' soon; for now, notice the qualitative law:
\emph{higher degree buys a wider window of trust.}

\section{The remainder: naming the gap honestly}

A picture that ``hugs'' a curve is suggestive but not yet mathematics. The honest
object is the \emph{remainder}, the exact difference between the function and its
polynomial:
\[
  R_n(x) = f(x) - P_n(x).
\]
This is not an approximation; it is an equation. It says $f(x) = P_n(x) + R_n(x)$
\emph{exactly}, for every $x$. All the error, all the discrepancy, has been gathered
up and given a name. The entire game of approximation is now the game of controlling
$R_n$: if we can show $R_n(x)$ is small, then $P_n(x)$ is a good approximation, with
a guarantee, not a hope.

\begin{happeningbox}
When you write $\cos x \approx 1 - \tfrac12 x^2$, the ``$\approx$'' is hiding a
precise statement. What is really there is an equation with a correction term:
\[
  \cos x = 1 - \tfrac12 x^2 + R_2(x),
\]
and Taylor's theorem is going to hand us a \emph{formula} for $R_2(x)$ --- not its
exact value (that would be asking too much), but a clean expression that lets us
\emph{bound} it. The difference between an engineer who trusts a small-angle
approximation and a gambler who hopes it works is exactly this remainder term. The
remainder is where approximation becomes a science.
\end{happeningbox}

So what should we expect of $R_n$? Two things, immediately. First, since $P_n$
matches $f$ and its first $n$ derivatives at $a$, the remainder and \emph{its} first
$n$ derivatives must all vanish at $a$:
\[
  R_n(a) = R_n'(a) = R_n''(a) = \cdots = R_n^{(n)}(a) = 0.
\]
A function whose first $n$ derivatives die at a point is extremely flat there ---
flatter than $(x-a)^n$. So $R_n(x)$ ought to behave like $(x-a)^{n+1}$ near $a$:
small, and getting small \emph{fast} as $x \to a$. That is the heart of why more
terms help. Let us turn this intuition into a theorem.

\section{Taylor's theorem with the Lagrange remainder}

Here is the destination, stated cleanly, before we earn it.

\begin{focusbox}
\textbf{Taylor's Theorem (Lagrange form of the remainder).} Suppose $f$ is $n+1$
times differentiable on an open interval containing $a$ and $x$. Then there is a
point $c$ strictly between $a$ and $x$ such that
\[
  f(x) = \underbrace{\sum_{k=0}^{n} \frac{f^{(k)}(a)}{k!}(x-a)^k}_{P_n(x)}
       \; + \; \underbrace{\frac{f^{(n+1)}(c)}{(n+1)!}\,(x-a)^{n+1}}_{R_n(x)}.
\]
The remainder $R_n(x)$ looks exactly like the \emph{next} Taylor term --- the
degree-$(n+1)$ term you would have written next --- except that the
$(n+1)$-st derivative is evaluated not at $a$ but at some mystery point $c$ in between.
\end{focusbox}

Look at how beautifully this generalizes the Mean Value Theorem. Set $n=0$. The
theorem says $f(x) = f(a) + f'(c)(x-a)$ for some $c$ between $a$ and $x$ --- which,
rearranged, is $\frac{f(x)-f(a)}{x-a} = f'(c)$, precisely the MVT. So Taylor's
theorem is the MVT with more derivatives. The single mysterious point $c$ where ``the
average is achieved'' in the MVT reappears here as the point where the
$(n+1)$-st derivative takes the value that makes the equation exactly balance.

\begin{intuitionbox}
The Lagrange remainder is the optimist's dream made rigorous. You were \emph{hoping}
the error would behave like the next term $\frac{f^{(n+1)}(a)}{(n+1)!}(x-a)^{n+1}$.
Taylor's theorem says: it does behave exactly like that next term --- you just have
to evaluate the derivative at a slightly shifted point $c$ instead of at $a$. And
since you usually do not need $c$ exactly, only a \emph{bound} on $f^{(n+1)}$ over
the interval, this shifted point costs you almost nothing in practice. You trade
``the exact error'' (impossible to know) for ``the error has this exact \emph{form}''
(immediately useful).
\end{intuitionbox}

\begin{historybox}
Brook Taylor published the polynomial expansion in 1715, though Newton, Gregory, and
the Indian Kerala school centuries earlier already knew the series for $\sin$,
$\cos$, and $\arctan$. What none of them had was the \emph{remainder} --- the precise
accounting of error. The early calculus treated these series as exact infinite sums
and computed with them by bold instinct, usually correctly. It was Joseph-Louis
Lagrange in 1797 who isolated the finite remainder term that bears his name, and
Cauchy and Weierstrass in the next century who made the whole apparatus rigorous by
grounding it in the limit and the Mean Value Theorem. The picture had been right for a
hundred years; the guarantee came later.
\end{historybox}

\subsection{Why it is true: the idea before the proof}

There are several routes to this theorem. The cleanest conceptually is a souped-up
version of Rolle's theorem, applied to an auxiliary function engineered so that its
vanishing \emph{is} the statement we want. Before the formalism, here is the idea in
one breath.

We want to compare $f(x)$ with $P_n(x)$. Fix the target point $x = b$ and the center
$a$. Build a helper function $g(t)$ of a new variable $t$ that measures, at each $t$,
how far $f$ is from \emph{its own Taylor polynomial centered at $t$} (evaluated at
$b$), and rig it so that $g$ equals $0$ at the end $t=b$. Then a single
mean-value step on $g$ --- this is the magic --- causes almost everything to
telescope away, leaving precisely the remainder formula. The telescoping is not
luck; the helper is designed to telescope.

\begin{whybox}
Why a helper function? Because the quantity we care about, $f(b) - P_n(b)$, is a
single number with no $t$ in it --- there is nothing to apply a mean-value argument
to. The trick, which runs through all of analysis, is to \emph{introduce a variable}:
turn the static number into the endpoint value of a moving quantity, so that the
machinery of derivatives (Rolle, MVT) has something to bite on. We promote the
fixed center $a$ to a sliding center $t$, watch the discrepancy change as $t$ slides
from $a$ to $b$, and let the Mean Value Theorem catch the instant where the change
conspires to reveal the remainder. Designing the right helper is the whole art;
verifying it is then mechanical.
\end{whybox}

\subsection{The proof}

\begin{theorem}[Taylor with Lagrange remainder]
\label{ch31:thm-taylor}
Let $f$ be a function such that $f, f', \dots, f^{(n)}$ are continuous on the closed
interval between $a$ and $b$, and $f^{(n+1)}$ exists on the open interval between
them. Then there is a point $c$ strictly between $a$ and $b$ with
\[
  f(b) = \sum_{k=0}^{n} \frac{f^{(k)}(a)}{k!}(b-a)^k
       + \frac{f^{(n+1)}(c)}{(n+1)!}(b-a)^{n+1}.
\]
\end{theorem}

\begin{proof}
Treat $a$, $b$, $n$ as fixed. Define a function of a sliding center $t$, living on
the interval between $a$ and $b$:
\[
  g(t) = f(b) - \sum_{k=0}^{n} \frac{f^{(k)}(t)}{k!}\,(b-t)^k.
\]
Read $g(t)$ out loud: it is the gap between the true value $f(b)$ and the Taylor
polynomial of $f$ \emph{centered at $t$}, evaluated at $b$. Two endpoint values are
immediate.

At $t = b$: every factor $(b-t)^k$ with $k \ge 1$ vanishes, and the $k=0$ term is
$f(b)$, so
\[
  g(b) = f(b) - f(b) = 0.
\]
At $t = a$: the sum is exactly the Taylor polynomial $P_n(b)$, so
\[
  g(a) = f(b) - P_n(b) = R_n(b),
\]
which is the very remainder we are chasing. Good: $g$ carries the answer at one end
and is zero at the other.

Now differentiate $g$ with respect to $t$. This is the step where the design pays
off. Each term of the sum is a product $\frac{f^{(k)}(t)}{k!}(b-t)^k$, so we use the
product rule:
\[
  \frac{\dd}{\dd t}\!\left[\frac{f^{(k)}(t)}{k!}(b-t)^k\right]
  = \frac{f^{(k+1)}(t)}{k!}(b-t)^k
    \;-\; \frac{f^{(k)}(t)}{k!}\,k\,(b-t)^{k-1}.
\]
The first piece is the derivative landing on $f^{(k)}$; the second is the derivative
landing on $(b-t)^k$, which brings down $k$ and lowers the power (the minus sign
comes from the chain rule on $b-t$). Now sum over $k = 0$ to $n$ and watch the
second pieces cancel the first pieces of the neighboring term. Writing the first-type
terms as $A_k = \frac{f^{(k+1)}(t)}{k!}(b-t)^k$ and the second-type terms as
$B_k = \frac{f^{(k)}(t)}{(k-1)!}(b-t)^{k-1}$ (note $\frac{k}{k!} = \frac{1}{(k-1)!}$),
we have $B_k = A_{k-1}$. So
\[
  g'(t) = -\sum_{k=0}^{n}\big(A_k - B_k\big)
        = -\sum_{k=0}^{n} A_k + \sum_{k=1}^{n} A_{k-1}
        = -A_n,
\]
because the $B$-sum telescopes against the $A$-sum, leaving only the very last
first-type term (the $B_0$ term is zero since it has a factor $k=0$). Therefore
\[
  g'(t) = -\frac{f^{(n+1)}(t)}{n!}\,(b-t)^n.
\]
Everything collapsed to a single clean term: this is exactly the telescoping the
helper was engineered to produce.

Now we apply a mean-value argument --- but cleverly, using a second helper so that we
get the strong $(n+1)!$ form rather than a weaker estimate. Introduce
\[
  h(t) = (b-t)^{n+1}.
\]
Both $g$ and $h$ are continuous on the closed interval and differentiable on the open
one, and $h'(t) = -(n+1)(b-t)^n$ never vanishes inside. The Cauchy (generalized) Mean
Value Theorem then guarantees a point $c$ strictly between $a$ and $b$ with
\[
  \frac{g(a) - g(b)}{h(a) - h(b)} = \frac{g'(c)}{h'(c)}.
\]
We computed $g(b) = 0$, $g(a) = R_n(b)$, $h(b) = 0$, and $h(a) = (b-a)^{n+1}$. The
right side is
\[
  \frac{g'(c)}{h'(c)}
  = \frac{-\frac{f^{(n+1)}(c)}{n!}(b-c)^n}{-(n+1)(b-c)^n}
  = \frac{f^{(n+1)}(c)}{(n+1)!}.
\]
The $(b-c)^n$ factors cancel --- another piece of the design --- and the $n!$ becomes
$(n+1)!$. Setting the two sides equal,
\[
  \frac{R_n(b)}{(b-a)^{n+1}} = \frac{f^{(n+1)}(c)}{(n+1)!},
\]
which rearranges to exactly
\[
  R_n(b) = \frac{f^{(n+1)}(c)}{(n+1)!}(b-a)^{n+1}.
\]
That is the claim.
\end{proof}

\begin{remark}
The case $n=0$ recovers the Mean Value Theorem itself, and the proof above, stripped
to $n=0$, \emph{is} the standard proof of the MVT via the Cauchy form. So Taylor's
theorem is not a different animal from the MVT; it is the same animal, fed more
derivatives. Everything in the chapter on the Mean Value Theorem was secretly
preparation for this.
\end{remark}

\section{Using the remainder to bound error}

The point $c$ is unknown, and that is fine, because we almost never want $c$ itself.
We want a guarantee. And a guarantee follows the instant we can say how big
$f^{(n+1)}$ can possibly be on the interval. This is the workhorse corollary --- the
single most useful fact in the whole chapter.

\begin{focusbox}
\textbf{The error bound.} If $\abs{f^{(n+1)}(t)} \le M$ for all $t$ between $a$ and
$x$, then
\[
  \abs{R_n(x)} \;\le\; \frac{M}{(n+1)!}\,\abs{x-a}^{n+1}.
\]
You do not need to find $c$. You only need a ceiling $M$ on the next derivative.
\end{focusbox}

This is the form you actually use. To approximate $f(x)$ by $P_n(x)$ and \emph{know}
how wrong you might be, you ask one question: how large can the $(n+1)$-st derivative
get on the way from $a$ to $x$? Whatever that ceiling $M$ is, the error cannot exceed
$\frac{M}{(n+1)!}\abs{x-a}^{n+1}$. Two forces fight in that bound: $\abs{x-a}^{n+1}$
grows if you reach far from $a$, but $(n+1)!$ in the denominator grows ferociously
with $n$. For functions whose derivatives stay bounded (like $\sin$, $\cos$, $e^x$ on
any fixed interval), the factorial wins decisively, and the error collapses to zero as
you take more terms.

\begin{labbox}
Approximate $\cos(0.5)$ using $P_2(x) = 1 - \tfrac12 x^2$ and bound the error
\emph{before} comparing to a calculator.

The estimate: $P_2(0.5) = 1 - \tfrac12(0.25) = 1 - 0.125 = 0.875$.

The bound: here $n=2$, so we need a ceiling on $\abs{f^{(3)}}$. For $f=\cos$,
$f'''(x) = \sin x$, and $\abs{\sin t} \le 1$ everywhere, so $M=1$. Then
\[
  \abs{R_2(0.5)} \le \frac{1}{3!}\,(0.5)^3 = \frac{0.125}{6} \approx 0.0208.
\]
So we may assert, with certainty, that the true $\cos(0.5)$ lies in
$[0.875 - 0.021,\; 0.875 + 0.021] = [0.854,\, 0.896]$.

The check: a calculator gives $\cos(0.5) = 0.87758\ldots$, and indeed
$\abs{0.87758 - 0.875} = 0.0026$, comfortably under our promised ceiling of
$0.0208$. The bound was honest and, as usual, a bit pessimistic --- which is exactly
what you want from a guarantee.
\end{labbox}

\begin{trapbox}
A bound is an upper limit, not a prediction. Beginners sometimes report ``the error
is $0.0208$'' when the theorem only promises ``the error is \emph{at most}
$0.0208$.'' The actual error here was $0.0026$, almost ten times smaller. The
Lagrange bound is deliberately conservative: it assumes the next derivative sits at
its worst-case value $M$ over the whole interval, which it usually does not. Never
quote the bound as the error; quote it as the worst case.
\end{trapbox}

\begin{labbox}
A harder, more satisfying one: how many terms of the Maclaurin series for $e^x$ do
we need to compute $e = e^1$ to within $10^{-6}$?

For $f(x) = e^x$ centered at $a=0$, every derivative is $e^x$ again, and on
$[0,1]$ we have $e^t \le e < 3$, so we may take $M = 3$. With $x=1$,
\[
  \abs{R_n(1)} \le \frac{3}{(n+1)!}\,\cdot 1^{n+1} = \frac{3}{(n+1)!}.
\]
We need $\frac{3}{(n+1)!} \le 10^{-6}$, i.e. $(n+1)! \ge 3\times 10^{6}$. Now just
march up the factorials: $9! = 362880$ (too small), $10! = 3628800 \ge 3\times10^6$.
So $(n+1) = 10$, meaning $n = 9$: the polynomial through the $x^9/9!$ term suffices.
Computing $\sum_{k=0}^{9} \frac{1}{k!}$ gives $2.7182815\ldots$ against the true
$e = 2.7182818\ldots$, an error of about $3\times 10^{-7}$ --- under our budget, as
promised. The factorial in the denominator is doing extraordinary work: ten terms and
we already have six digits.
\end{labbox}

\begin{applicationbox}
This is precisely how a calculator or a math library computes $e^x$, $\sin x$, and
the rest. It does not have a transcendental oracle inside; it adds up a few Taylor
terms (after first reducing the argument to a small range where convergence is fast,
e.g. using $\sin(x + 2\pi) = \sin x$ and angle-halving), and the Lagrange remainder
is what certifies that the result is correct to the last displayed digit. Every time
you press the cosine button, a remainder estimate guaranteed the answer before it ever
reached the screen.
\end{applicationbox}

\section{When does the series converge to the function?}

Now the bigger dream. Suppose $f$ has \emph{all} derivatives. We can keep building
Taylor polynomials of every degree, and we may ask whether the infinite series
\[
  \sum_{k=0}^{\infty} \frac{f^{(k)}(a)}{k!}(x-a)^k
\]
adds up to $f(x)$. The series equals $f(x)$ at a given point $x$ exactly when the
remainder vanishes in the limit:
\[
  \boxed{\;\lim_{n\to\infty} R_n(x) = 0\;}
\]
This is the whole truth of convergence, and it is worth saying slowly. The Taylor
series converges \emph{to $f$} at $x$ if and only if the leftover $R_n(x)$ shrinks to
zero as you include more terms. Not ``if the series converges'' --- the series can
converge to the \emph{wrong} thing, as we will see --- but if and only if the gap
between $f$ and the polynomial closes.

\begin{whybox}
Why is $R_n \to 0$ the exact condition, and not something about the series
``converging''? Because $f(x) = P_n(x) + R_n(x)$ is an identity for every $n$.
Rearranged, $R_n(x) = f(x) - P_n(x)$. As $n \to \infty$, $P_n(x)$ is the partial sum
of the Taylor series. So $P_n(x) \to f(x)$ --- the series summing to the function ---
is literally the same statement as $R_n(x) \to 0$. There is no daylight between them.
This is why all the labor went into the remainder: it is the only thing standing
between the polynomial and the function.
\end{whybox}

For the friendly functions, the Lagrange bound closes the case in one line. Take
$\sin x$ centered at $0$. Every derivative is one of $\pm\sin x, \pm\cos x$, so
$\abs{f^{(n+1)}(t)} \le 1$ for all $t$, giving
\[
  \abs{R_n(x)} \le \frac{1}{(n+1)!}\,\abs{x}^{n+1}.
\]
For any fixed $x$, no matter how large, $\frac{\abs{x}^{n+1}}{(n+1)!} \to 0$ as
$n\to\infty$ --- the factorial outruns any fixed power. Therefore the Maclaurin
series for $\sin x$ converges to $\sin x$ for \emph{every} real $x$. The same
argument, with $M = e^{\abs{x}}$, works for $e^x$. These functions are reconstructed
perfectly, everywhere, from their single point of data at the origin.

\begin{figure}[H]
\centering
\begin{tikzpicture}
\begin{axis}[
  width=13cm, height=7cm,
  axis lines=middle,
  xlabel={$x$}, ylabel={$y$},
  xmin=-7, xmax=7, ymin=-2.2, ymax=2.2,
  samples=200, domain=-6.8:6.8,
  legend pos=north west, legend cell align=left,
  restrict y to domain=-2.4:2.4,
  clip=true,
]
\addplot[curveblue, very thick] {sin(deg(x))};
\addlegendentry{$\sin x$}
\addplot[gray, thick, dashed] {x - x^3/6};
\addlegendentry{degree $3$}
\addplot[quotientorange, thick] {x - x^3/6 + x^5/120 - x^7/5040};
\addlegendentry{degree $7$}
\addplot[tangentred, thick] {x - x^3/6 + x^5/120 - x^7/5040 + x^9/362880 - x^11/39916800};
\addlegendentry{degree $11$}
\end{axis}
\end{tikzpicture}
\caption{Taylor polynomials of $\sin x$ at $0$ closing in on the curve. Each higher
degree extends the interval over which the polynomial is faithful before it finally
shoots off to infinity. In the limit, the series captures $\sin x$ on the entire real
line.}
\label{ch31:fig-sin-taylor}
\end{figure}

\subsection{Where convergence stops: the geometric warning}

Not every function is captured on the whole line. Consider the homely
\[
  f(x) = \frac{1}{1+x^2}.
\]
This function is perfectly smooth everywhere on $\R$ --- no breaks, no corners, no
vertical asymptotes, infinitely differentiable. Its Maclaurin series, obtained from
the geometric series $\frac{1}{1+u} = 1 - u + u^2 - \cdots$ with $u = x^2$, is
\[
  1 - x^2 + x^4 - x^6 + \cdots,
\]
which converges only for $\abs{x} < 1$ and diverges for $\abs{x} > 1$. Why $1$? There
is nothing visibly wrong with the function at $x = 1$. The graph sails smoothly
through $(1, \tfrac12)$ without a hint of trouble.

\begin{figure}[H]
\centering
\begin{tikzpicture}
\begin{axis}[
  width=12cm, height=6.5cm,
  axis lines=middle,
  xlabel={$x$}, ylabel={$y$},
  xmin=-2.2, xmax=2.2, ymin=-0.6, ymax=1.5,
  samples=160, domain=-2.1:2.1,
  legend pos=north west, legend cell align=left,
  restrict y to domain=-0.8:1.7,
  clip=true,
]
\addplot[curveblue, very thick] {1/(1+x^2)};
\addlegendentry{$\frac{1}{1+x^2}$}
\addplot[quotientorange, thick, dashed, domain=-1.4:1.4] {1 - x^2 + x^4 - x^6 + x^8};
\addlegendentry{partial sum (deg $8$)}
\addplot[only marks, mark=*, mark size=1.4pt, black] coordinates {(1,0.5) (-1,0.5)};
\draw[trapred, thick, dotted] (axis cs:1,-0.6) -- (axis cs:1,1.5);
\draw[trapred, thick, dotted] (axis cs:-1,-0.6) -- (axis cs:-1,1.5);
\end{axis}
\end{tikzpicture}
\caption{The function $\frac{1}{1+x^2}$ is smooth everywhere, yet its Maclaurin series
$1-x^2+x^4-\cdots$ tracks it only on $(-1,1)$ (dotted lines). Inside that interval the
partial sums close in; just outside, they fly apart, even though the function itself is
perfectly well behaved there. The wall at $\abs{x}=1$ is set by a singularity hiding off
the real axis.}
\label{ch31:fig-geometric}
\end{figure}

\begin{historybox}
The mystery dissolves the moment you allow complex numbers, and this was one of the
great clarifying insights of nineteenth-century analysis (Cauchy, then
Weierstrass and Riemann). The function $\frac{1}{1+x^2}$, viewed over the complex
plane, \emph{blows up} at $x = \pm i$, where $1 + x^2 = 0$. Those poles sit at
distance $1$ from the origin. A power series centered at $0$ converges inside the
largest disk around $0$ that reaches no singularity --- and the nearest singularities,
invisible on the real line, are exactly one unit away in the imaginary direction. The
radius of convergence $1$ was being dictated by points you could not see. The early
calculus of Newton and Euler used such series freely and often correctly by
instinct; it took the complex viewpoint to explain \emph{why} the instinct sometimes
failed.
\end{historybox}

\begin{connectionbox}
This is the bridge to complex analysis. A real power series centered at $a$ converges
on an interval $(a-R, a+R)$ (the dotted wall in Figure~\ref{ch31:fig-geometric});
the number $R$ is the \emph{radius of convergence}, and
it equals the distance from $a$ to the nearest point --- real \emph{or complex} ---
where the function misbehaves. The interval on the real line is just the diameter of a
hidden disk in the complex plane. Once you know this, ``why does this series stop
converging at $x=1$ when the function looks fine?'' is never mysterious again: look
for the nearest singularity off the real axis.
\end{connectionbox}

\subsection{The shocking counterexample: smooth but not analytic}

There is a worse failure, and it is worth meeting because it draws the sharp line
between ``has a Taylor series'' and ``equals its Taylor series.'' Consider
\[
  f(x) =
  \begin{cases}
    e^{-1/x^2}, & x \ne 0, \\
    0, & x = 0.
  \end{cases}
\]
This function is infinitely differentiable everywhere, including at $0$. But it is so
extraordinarily flat at the origin that one can show
\[
  f(0) = f'(0) = f''(0) = \cdots = f^{(n)}(0) = \cdots = 0
\]
--- \emph{every} derivative at $0$ is zero. The exponential $e^{-1/x^2}$ decays to
zero faster than any power of $x$ can detect, so the difference quotients defining
each derivative all collapse to $0$.

Now build the Maclaurin series. Every coefficient $\frac{f^{(k)}(0)}{k!}$ is $0$, so
the Taylor series is
\[
  0 + 0\cdot x + 0\cdot x^2 + \cdots = 0.
\]
The series converges beautifully --- to the zero function. But $f(x)$ is \emph{not}
zero for $x \ne 0$. The Taylor series converges everywhere and yet equals $f$ only at
the single point $x=0$. Here $R_n(x) = f(x) - 0 = f(x)$ does \emph{not} go to zero;
it just sits there equal to the whole function. The remainder never shrinks because
there is nothing in the polynomial to shrink toward $f$.

\begin{figure}[H]
\centering
\begin{tikzpicture}
\begin{axis}[
  width=11cm, height=6cm,
  axis lines=middle,
  xlabel={$x$}, ylabel={$y$},
  xmin=-2.6, xmax=2.6, ymin=-0.15, ymax=1.05,
  samples=200, domain=-2.5:2.5,
  clip=true,
]
\addplot[curveblue, very thick] {exp(-1/(x^2+0.0000001))};
\addplot[tangentred, thick, dashed, domain=-2.5:2.5] {0};
\addplot[only marks, mark=*, mark size=1.5pt, black] coordinates {(0,0)};
\node[curveblue] at (axis cs:1.9,0.85) {$e^{-1/x^2}$};
\node[tangentred] at (axis cs:-1.4,0.08) {Taylor series $\equiv 0$};
\end{axis}
\end{tikzpicture}
\caption{The function $e^{-1/x^2}$ (extended by $0$ at the origin) is infinitely
flat at $x=0$: every derivative there is zero. Its Taylor series is identically zero
(red), and matches the function only at the single point $x=0$. Smoothness does not
guarantee that a function equals its Taylor series.}
\label{ch31:fig-flat}
\end{figure}

\begin{focusbox}
\textbf{Analytic functions.} A function is \emph{analytic} at $a$ if its Taylor
series centered at $a$ converges to the function on some interval around $a$ --- that
is, if $R_n(x) \to 0$ there. Most familiar functions ($e^x$, $\sin$, $\cos$,
polynomials, $\ln$, $1/(1-x)$ away from their singularities) are analytic. The
example $e^{-1/x^2}$ shows that \emph{infinitely differentiable is strictly weaker
than analytic}: being smooth is necessary but not sufficient. Equality with the
Taylor series is an extra, special grace, and it is exactly the grace that
$R_n \to 0$ certifies.
\end{focusbox}

\begin{aestheticbox}
There is something humbling and beautiful here. We began the chapter dreaming that one
point, fully known, determines the whole function. For analytic functions the dream
is literally true: $e^x$ is the unique analytic function with $f(0)=1$ and $f'=f$, and
its values everywhere are encoded in its behavior at the origin. But the dream has an
edge. The function $e^{-1/x^2}$ shows that a smooth function can keep secrets ---
information that simply is not present in any derivative at $0$. Analyticity is the
property of having \emph{no secrets}: every value is foretold by the local data. That
two such different worlds live side by side among the smooth functions is one of the
quietly profound facts of analysis.
\end{aestheticbox}

\section{Physics lives in the first two terms}

A great deal of physics is secretly the statement ``keep the first Taylor term that
does not vanish, throw away the rest, and the world becomes solvable.'' Approximation
is not a compromise physicists tolerate; it is often where the insight is.

\subsection{The small-angle approximation}

The pendulum equation, the bending of light, the behavior of a vibrating string ---
all of them simplify enormously the moment you replace $\sin\theta$ by its first
Taylor term. From the series
\[
  \sin\theta = \theta - \frac{\theta^3}{6} + \frac{\theta^5}{120} - \cdots,
\]
for small $\theta$ we keep just the leading term:
\[
  \sin\theta \approx \theta \qquad (\theta \text{ in radians}).
\]
The remainder bound tells you exactly when this is safe. The first dropped term is
$\theta^3/6$, and (using $\abs{f^{(3)}} \le 1$) the Lagrange remainder for the
degree-$1$ approximation is at most $\frac{\abs{\theta}^3}{6}$. At
$\theta = 0.1$ rad (about $5.7^\circ$), the error is below
$\frac{(0.1)^3}{6} \approx 1.7\times 10^{-4}$ --- a relative error well under a tenth
of a percent. This single approximation is why the pendulum has a period
$T = 2\pi\sqrt{L/g}$ \emph{independent of amplitude}: replacing $\sin\theta$ by
$\theta$ turns the true (hard) pendulum equation into the simple harmonic oscillator,
whose period does not depend on how far it swings. The famous result is an artifact of
truncating a Taylor series at the first term.

\begin{figure}[H]
\centering
\begin{tikzpicture}
\begin{axis}[
  width=11cm, height=6.5cm,
  axis lines=middle,
  xlabel={$\theta$ (radians)}, ylabel={},
  xmin=-1.6, xmax=1.6, ymin=-1.4, ymax=1.4,
  samples=160, domain=-1.55:1.55,
  legend pos=south east, legend cell align=left,
  clip=true,
]
\addplot[curveblue, very thick] {sin(deg(x))};
\addlegendentry{$\sin\theta$}
\addplot[tangentred, thick, dashed] {x};
\addlegendentry{$\theta$ (small-angle)}
\end{axis}
\end{tikzpicture}
\caption{The small-angle approximation $\sin\theta \approx \theta$. Near the origin
the line $y=\theta$ and the curve $y=\sin\theta$ are nearly indistinguishable; the gap
grows like $\theta^3/6$, so it stays tiny while $\theta$ is small.}
\label{ch31:fig-small-angle}
\end{figure}

\begin{applicationbox}
The same ``first nonvanishing term'' logic runs through physics:
\begin{itemize}
\item \textbf{Relativistic energy.} The exact kinetic energy is
$E = mc^2\big(\tfrac{1}{\sqrt{1-v^2/c^2}} - 1\big)$. Expanding the square root for
small $v/c$ gives $E \approx \tfrac12 mv^2 + \tfrac{3}{8}\frac{mv^4}{c^2} + \cdots$.
The first term is Newton's $\tfrac12 mv^2$ --- so classical kinetic energy is just the
leading Taylor term of the relativistic one. The next term is the first relativistic
correction.
\item \textbf{Potential wells.} Near a stable equilibrium $x_0$, any smooth potential
$V(x)$ has $V'(x_0)=0$, so its Taylor expansion begins
$V(x) \approx V(x_0) + \tfrac12 V''(x_0)(x-x_0)^2$. The quadratic term is a spring:
\emph{every} stable system, looked at closely enough, oscillates like a harmonic
oscillator. This is why $\tfrac12 k x^2$ appears everywhere in physics.
\item \textbf{Optics and pendulums.} $\tan\theta \approx \theta$ and
$\cos\theta \approx 1 - \tfrac12\theta^2$ linearize lens equations and give the
leading correction to a pendulum's period.
\end{itemize}
In each case the Taylor remainder is the auditor that tells you when the
approximation may be trusted and when the next term must be restored.
\end{applicationbox}

\begin{walkbox}
Walk and turn this over: why is it the \emph{second} derivative that governs the
behavior of a system near equilibrium, never the first? (Hint: equilibrium means the
first derivative of the potential is zero, so the first surviving Taylor term is the
quadratic one --- and a quadratic potential is exactly a spring. Stability is
$V''>0$; instability is $V''<0$.) What would a system look like if $V''$ also
vanished and the leading term were quartic?
\end{walkbox}

\section{Worked examples, start to finish}

\begin{example}[Estimating $\sqrt{e}$ with a guarantee]
\label{ch31:ex-sqrt-e}
Approximate $\sqrt{e} = e^{0.5}$ using the degree-$3$ Maclaurin polynomial of $e^x$
and bound the error.

The polynomial is $P_3(x) = 1 + x + \tfrac{x^2}{2} + \tfrac{x^3}{6}$. At $x=0.5$,
\[
  P_3(0.5) = 1 + 0.5 + \tfrac{0.25}{2} + \tfrac{0.125}{6}
           = 1 + 0.5 + 0.125 + 0.0208\overline{3} = 1.6458\overline{3}.
\]
For the error, $f^{(4)}(x) = e^x$ and on $[0,0.5]$ we have
$e^t \le e^{0.5} < 1.65$, so take $M = 1.65$. Then
\[
  \abs{R_3(0.5)} \le \frac{1.65}{4!}(0.5)^4
  = \frac{1.65}{24}\cdot 0.0625 \approx 0.0043.
\]
So $\sqrt{e} \in [1.6458 - 0.0043,\; 1.6458 + 0.0043] = [1.6415,\, 1.6501]$. The true
value is $\sqrt{e} = 1.64872\ldots$, inside the interval, with actual error
$\approx 0.0029$ --- under the promised $0.0043$.
\end{example}

\begin{example}[A Taylor expansion not centered at $0$]
\label{ch31:ex-ln}
Approximate $\ln(1.2)$ by expanding $\ln x$ about $a = 1$.

The derivatives of $f(x) = \ln x$ are $f'(x) = x^{-1}$, $f''(x) = -x^{-2}$,
$f'''(x) = 2x^{-3}$, so at $a=1$: $f(1)=0$, $f'(1)=1$, $f''(1)=-1$, $f'''(1)=2$. The
degree-$2$ polynomial is
\[
  P_2(x) = 0 + 1\cdot(x-1) + \frac{-1}{2}(x-1)^2 = (x-1) - \tfrac12(x-1)^2.
\]
At $x = 1.2$, with $(x-1) = 0.2$:
\[
  P_2(1.2) = 0.2 - \tfrac12(0.04) = 0.2 - 0.02 = 0.18.
\]
For the error, $f'''(x) = 2x^{-3}$ and on $[1, 1.2]$ this is largest at $x=1$, where
it equals $2$, so $M = 2$. Then
\[
  \abs{R_2(1.2)} \le \frac{2}{3!}(0.2)^3 = \frac{2}{6}(0.008) \approx 0.00267.
\]
The true value $\ln(1.2) = 0.18232\ldots$ differs from $0.18$ by $0.00232$, inside
the bound. Notice we expanded about $1$, not $0$, because $\ln$ is undefined at $0$ ---
you center where the function is friendly and where you know the derivatives.
\end{example}

\begin{example}[Reading off a limit from the first surviving term]
\label{ch31:ex-limit}
Evaluate $\displaystyle\lim_{x\to 0}\frac{x - \sin x}{x^3}$ using Taylor expansion.

Expand the numerator: $\sin x = x - \tfrac{x^3}{6} + \tfrac{x^5}{120} - \cdots$, so
\[
  x - \sin x = \frac{x^3}{6} - \frac{x^5}{120} + \cdots.
\]
Divide by $x^3$:
\[
  \frac{x - \sin x}{x^3} = \frac{1}{6} - \frac{x^2}{120} + \cdots
  \;\xrightarrow[x\to 0]{}\; \frac{1}{6}.
\]
The limit is $\tfrac16$. This is the same answer L'Hopital's rule gives after three
differentiations, but Taylor expansion makes the mechanism transparent: the
numerator's first nonzero term is $x^3/6$, and dividing by $x^3$ exposes its
coefficient. Taylor series turn many delicate limits into bookkeeping of leading
terms.
\end{example}

\begin{connectionbox}
Example~\ref{ch31:ex-limit} is no accident. Whenever you face a $0/0$ limit, expanding
numerator and denominator in Taylor series and comparing their lowest surviving powers
resolves it instantly --- and explains \emph{why} L'Hopital's rule works at all.
L'Hopital differentiates to peel away vanishing orders one at a time; Taylor expansion
shows all the orders at once. They are two views of the same fact: near a zero, a
smooth function looks like its lowest nonvanishing Taylor term.
\end{connectionbox}

\section{What to carry away}

\begin{itemize}
\item The degree-$n$ \textbf{Taylor polynomial} $P_n(x) = \sum_{k=0}^n
\frac{f^{(k)}(a)}{k!}(x-a)^k$ is the unique polynomial matching $f$ and its first $n$
derivatives at $a$. The factorials exactly undo the factorials that repeated
differentiation produces.
\item \textbf{Taylor's theorem} makes the approximation exact:
$f(x) = P_n(x) + \frac{f^{(n+1)}(c)}{(n+1)!}(x-a)^{n+1}$ for some $c$ between $a$ and
$x$. The remainder is the next term with its derivative evaluated at a hidden point.
It generalizes the Mean Value Theorem (the case $n=0$) and is proved the same way.
\item The \textbf{error bound}
$\abs{R_n(x)} \le \frac{M}{(n+1)!}\abs{x-a}^{n+1}$, where $M$ bounds $f^{(n+1)}$,
is the practical workhorse. It needs only a ceiling on the next derivative, never the
mystery point $c$.
\item The Taylor series equals $f(x)$ \textbf{exactly when $R_n(x)\to 0$}. For
$e^x, \sin, \cos$ the factorial crushes the remainder everywhere. For $\frac{1}{1+x^2}$
a hidden complex singularity caps the radius of convergence. For $e^{-1/x^2}$ the
series converges to the wrong function entirely: smooth does not imply analytic.
\item \textbf{Analytic} functions are those that equal their Taylor series --- the
ones with no local secrets. Physics lives largely in their first surviving terms:
$\sin\theta\approx\theta$, $V\approx \tfrac12 V'' (x-x_0)^2$, $E\approx\tfrac12 mv^2$.
\end{itemize}

\section{Exercises}

\begin{exercise}
Find the degree-$3$ Maclaurin polynomial of $f(x) = e^{x}\cos x$ by computing
$f(0), f'(0), f''(0), f'''(0)$ directly. (You should get
$1 + x - \tfrac{x^3}{3}$.)
\end{exercise}

\begin{exercise}
Use the degree-$2$ Taylor polynomial of $\sqrt{x}$ centered at $a = 4$ to estimate
$\sqrt{4.2}$. Then use the Lagrange remainder to bound the error, and compare with a
calculator value.
\end{exercise}

\begin{exercise}
The small-angle approximation $\sin\theta \approx \theta$ is used for
$\theta = 10^\circ$. Convert to radians, and use the remainder bound
$\abs{R_1(\theta)} \le \abs{\theta}^3/6$ to find the maximum possible error. Is the
approximation good to three decimal places?
\end{exercise}

\begin{exercise}
How many terms of the Maclaurin series for $\cos x$ are needed to compute
$\cos(1)$ with error less than $10^{-8}$? Justify with the Lagrange bound, then state
the polynomial you would use.
\end{exercise}

\begin{exercise}
Use Taylor expansions to evaluate
$\displaystyle \lim_{x\to 0}\frac{e^{x} - 1 - x}{x^2}$. Identify the first surviving
term of the numerator and read off the answer.
\end{exercise}

\begin{exercise}[Conceptual]
A classmate says, ``If a function is infinitely differentiable, its Taylor series must
converge to it.'' Give the counterexample from this chapter and explain in one or two
sentences exactly which step of that reasoning fails.
\end{exercise}

\begin{exercise}[Conceptual]
The Maclaurin series of $\frac{1}{1-x}$ is $1 + x + x^2 + \cdots$, with radius of
convergence $1$. Explain why $1$ is the radius, both from the real-line behavior of
the function and from the location of its singularity. Then predict the radius of
convergence of the Maclaurin series of $\frac{1}{4+x^2}$ without computing the series,
and justify your prediction.
\end{exercise}

\begin{exercise}[Proof]
Prove the integral form of the remainder: if $f^{(n+1)}$ is continuous on the interval
between $a$ and $x$, then
\[
  R_n(x) = \frac{1}{n!}\int_a^{x} f^{(n+1)}(t)\,(x-t)^n\,\dd t.
\]
(Hint: start from $f(x) - f(a) = \int_a^x f'(t)\,\dd t$ and integrate by parts
repeatedly, choosing the antiderivative of $1$ to be $-(x-t)$ so the powers build up.
This is the cleanest remainder of all, and the Lagrange form follows from it by the
Mean Value Theorem for integrals.)
\end{exercise}

\begin{exercise}[Proof]
Show that for $f(x) = e^{-1/x^2}$ (with $f(0)=0$), the first derivative at $0$ is
$f'(0) = 0$. (Compute $\lim_{h\to 0}\frac{e^{-1/h^2}}{h}$ by substituting
$u = 1/h^2 \to \infty$ and comparing the growth of $e^{u}$ against $\sqrt{u}$.) Explain
in words why the same kind of argument forces \emph{every} derivative at $0$ to vanish.
\end{exercise}

\begin{exercise}
Expand $f(x) = \frac{1}{\sqrt{1+x}}$ to degree $2$ about $a=0$ to derive the
approximation $\frac{1}{\sqrt{1+x}} \approx 1 - \tfrac12 x + \tfrac38 x^2$. Use it to
estimate $\frac{1}{\sqrt{1.1}}$, bound the error, and explain how this is the Taylor
origin of the relativistic correction term mentioned in the chapter.
\end{exercise}
